% =====================================================================================
%  Shared look for the handout and the speaker copy.  Compiled with XeLaTeX.
%  Edit colours, fonts and the code-block style here.
% =====================================================================================
\usepackage{geometry}
\usepackage{xcolor}
\usepackage{fontspec}
\usepackage{fancyhdr}
\usepackage{graphicx}
\usepackage{tikz}
\usetikzlibrary{calc}
\usepackage{listings}

% ------------------------------------------------------------------ palette
\definecolor{pPink}{HTML}{FAD2E7}
\definecolor{pMint}{HTML}{D2FAD3}
\definecolor{pLav}{HTML}{DFD2FA}
\definecolor{pPeach}{HTML}{F5DAC9}
\definecolor{ink}{HTML}{5E5870}
\definecolor{acc}{HTML}{8E4F86}
\definecolor{teal}{HTML}{3F8A8C}
\colorlet{colbg}{pPeach!62!white}        % shaded speaker-script column
\colorlet{colrule}{pPeach!55!ink!45}
\definecolor{codebg}{HTML}{F4EFFC}       % code blocks
\definecolor{codecomment}{HTML}{7D7793}
\definecolor{codenum}{HTML}{A59FB8}
\definecolor{codenumber}{HTML}{B5651D}

% ------------------------------------------------------------------ fonts (names come from config.yml)
% Each font falls back to Latin Modern when it is not installed, and logs a note that build.py reports.
\newcommand{\TryFont}[3]{\IfFontExistsTF{#1}{#2}{\typeout{talk: font '#1' not installed}#3}}
\TryFont{\FontBody}{\setmainfont{\FontBody}}{\setmainfont{Latin Modern Roman}}
\TryFont{\FontHeading}{\setsansfont{\FontHeading}}{\setsansfont{Latin Modern Sans}}
\TryFont{\FontCode}{\setmonofont[Scale=MatchLowercase]{\FontCode}}{\setmonofont[Scale=MatchLowercase]{Latin Modern Mono}}
% the speaker script gets its own family so it never looks like the handout text
\TryFont{\FontScript}{\newfontfamily\ScriptText{\FontScript}}{\newfontfamily\ScriptText{Latin Modern Roman}}

\setlength{\parindent}{0pt}
\setlength{\parskip}{0.55em}
\linespread{1.07}
\setlength{\leftmargini}{1.4em}
\setlength{\leftmarginii}{1.2em}
\providecommand{\tightlist}{\setlength{\itemsep}{0pt}\setlength{\parskip}{0pt}}
\providecommand{\passthrough}[1]{#1}
\providecommand{\pandocbounded}[1]{#1}

% ------------------------------------------------------------------ code blocks
%  \codesize is redefined by the speaker column so code shrinks together with the script text.
\newcommand{\codesize}{\fontsize{8.6}{10.8}\selectfont}

\lstdefinestyle{talk}{
  basicstyle=\ttfamily\codesize\color{ink},
  keywordstyle=\color{acc}\bfseries,
  keywordstyle={[2]\color{teal}\bfseries},
  stringstyle=\color{teal},
  commentstyle=\color{codecomment}\itshape,
  identifierstyle=\color{ink},
  numberstyle=\ttfamily\fontsize{6.5}{8}\selectfont\color{codenum},
  numbers=left, numbersep=6pt, stepnumber=1,
  backgroundcolor=\color{codebg},
  frame=l, framerule=2.6pt, rulecolor=\color{acc},
  xleftmargin=27pt, framexleftmargin=24pt, framexrightmargin=0pt,
  framextopmargin=4pt, framexbottommargin=4pt,
  aboveskip=0.9em, belowskip=0.7em,
  breaklines=true, breakatwhitespace=false, breakindent=1.2em,
  postbreak=\mbox{\textcolor{acc}{\scriptsize$\hookrightarrow$\,}},
  columns=fullflexible, keepspaces=true, tabsize=4, showstringspaces=false,
  extendedchars=true, escapeinside={(*@}{@*)},
}
\lstset{style=talk}

% languages listings does not know
\lstdefinelanguage{json}{
  morestring=[b]", sensitive=true, morekeywords={true,false,null},
}
\lstdefinelanguage{yaml}{
  morestring=[b]", morestring=[b]', morecomment=[l]{\#}, sensitive=true,
  morekeywords={true,false,null,yes,no},
}

% ------------------------------------------------------------------ header / footer
\newcommand{\CurSlide}{0}
\newcommand{\DocKind}{}
\pagestyle{fancy}
\fancyhf{}
\renewcommand{\headrulewidth}{0pt}
\fancyfoot[L]{\sffamily\scriptsize\color{ink}\TalkTitle\ \textperiodcentered\ \DocKind}
\fancyfoot[R]{\sffamily\scriptsize\color{ink}Slide \CurSlide\ / \TalkNumSlides}

% ------------------------------------------------------------------ pieces
\newcommand{\SlideFrame}[1]{%
  \begingroup\setlength{\fboxsep}{0pt}\setlength{\fboxrule}{0.5pt}%
  \fcolorbox{ink!35}{white}{\includegraphics[width=\dimexpr\linewidth-2\fboxrule\relax]{#1}}%
  \endgroup}

\newcommand{\RuledLines}[1]{% fills the remaining vertical space with ruled lines, #1 = spacing
  \leaders\vbox to #1{\vss\hbox to \linewidth{\color{ink!35}\leaders\hrule height 0.4pt\hfill}}\vfill}

\newcommand{\CoverPage}{%
  \thispagestyle{empty}%
  \begin{tikzpicture}[remember picture,overlay]
    \fill[pLav] (current page.north west) rectangle ([yshift=-72mm]current page.north east);
    \fill[pPink] ([yshift=-72mm]current page.north west) rectangle ([yshift=-76mm]current page.north east);
    \fill[pMint] ([yshift=-76mm]current page.north west) rectangle ([yshift=-79mm]current page.north east);
  \end{tikzpicture}%
  \vspace*{22mm}%
  {\sffamily\bfseries\fontsize{30}{34}\selectfont\color{ink}\TalkTitle\par}%
  \vspace{4mm}%
  {\sffamily\fontsize{15}{19}\selectfont\color{ink}\TalkSubtitle\par}%
  \vspace{62mm}%
  {\sffamily\large\color{acc}\bfseries\DocKind\par}%
  \vspace{3mm}%
  {\sffamily\color{ink}\TalkSpeaker\par\TalkEvent\par\TalkDate\par}%
  \clearpage}

\newcommand{\HandoutPage}[2]{%
  \renewcommand{\CurSlide}{#2}%
  \SlideFrame{build/slides/slide-#1.png}\par
  \vspace{4mm}%
  \input{build/tex/handout-#1.tex}\par
  \ifBlankLines\vspace{2mm}\RuledLines{8.5mm}\fi
  \clearpage}

% ------------------------------------------------------------------ speaker column
%  The script is typeset in a box exactly as wide as the column. It is tried at 11 pt and shrinks
%  step by step (to 8.5 pt) until it leaves \MinNotes of ruled "Notes" space. If it still does not
%  fit, it is set at 9.5 pt and continues in the column on extra pages (\vsplit), so nothing is lost.
\newlength{\ColInnerW}\newlength{\ColInnerH}\newlength{\ScriptMaxH}
\AtBeginDocument{%
  \setlength{\ColInnerW}{\dimexpr\ColW-10mm\relax}%
  \setlength{\ColInnerH}{\dimexpr\paperheight-22mm\relax}%
  % column height minus header block (12 mm) and the NOTES label/gap (9 mm) minus the minimum notes lines
  \setlength{\ScriptMaxH}{\dimexpr\ColInnerH-21mm-\MinNotes\relax}}

\newsavebox{\ScriptMeasure}
\newbox\ScriptRest
\newbox\ScriptPart
\newif\ifScriptOver
\newcommand{\ScriptFile}{}

% #1 font size, #2 leading, #3 code size, #4 code leading
\newcommand{\ScriptBody}[4]{{%
  \ScriptText\fontsize{#1}{#2}\selectfont
  \renewcommand{\codesize}{\fontsize{#3}{#4}\selectfont}%
  \raggedright\setlength{\parskip}{0.55em}\setlength{\parindent}{0pt}%
  \setlength{\leftmargini}{1.2em}%
  \input{\ScriptFile}\par}}

% a plain vbox (not a minipage) so that \vsplit can break it between lines
\makeatletter
\newcommand{\ScriptBox}[4]{\vbox{\hsize=\ColInnerW \@parboxrestore \ScriptBody{#1}{#2}{#3}{#4}}}
\makeatother

% Try a size; if the script does not leave room for the notes, try the next smaller one (#5).
\newcommand{\TryFit}[5]{%
  \setbox\ScriptMeasure\ScriptBox{#1}{#2}{#3}{#4}%
  \ifdim\dimexpr\ht\ScriptMeasure+\dp\ScriptMeasure\relax>\ScriptMaxH
    #5%
  \else
    \def\ChosenSize{{#1}{#2}{#3}{#4}}%
  \fi}

\newcommand{\FitScript}{%
  \ScriptOverfalse
  \TryFit{11}{14.6}{9.5}{12}{%
  \TryFit{10.5}{13.7}{9}{11.4}{%
  \TryFit{10}{13}{8.6}{10.8}{%
  \TryFit{9.5}{12.3}{8.2}{10.3}{%
  \TryFit{9}{11.7}{7.8}{9.8}{%
  \TryFit{8.5}{11}{7.4}{9.3}{%
    \ScriptOvertrue
    \def\ChosenSize{{9.5}{12.3}{8.2}{10.3}}}}}}}}}

\newcommand{\ColumnLabel}[1]{{\sffamily\bfseries\scriptsize\color{ink}#1}}

% One column page: header, a chunk of script taken from \ScriptRest, then ruled notes lines to the bottom.
\newcommand{\ColumnPage}[2]{%
  \ifdim\dimexpr\ht\ScriptRest+\dp\ScriptRest\relax>\ScriptMaxH
    \splittopskip=0pt \splitmaxdepth=\maxdimen \vbadness=10000
    \setbox\ScriptPart\vsplit\ScriptRest to \ScriptMaxH
  \else
    \setbox\ScriptPart\box\ScriptRest
  \fi
  \vbox to \ColInnerH{\hsize=\ColInnerW\linewidth=\ColInnerW\parindent=0pt\parskip=0pt\normalfont
    \hbox to \hsize{\ColumnLabel{#1}\hfill\ColumnLabel{SLIDE #2}}%
    \vskip 1.6mm
    \hbox to \hsize{\color{colrule}\leaders\hrule height 0.6pt\hfill}%
    \vskip 3.5mm
    \box\ScriptPart
    \vskip 5mm
    \ColumnLabel{NOTES}\par
    \vskip 0.5mm
    \leaders\vbox to 8mm{\vss\hbox to \ColInnerW{\color{colrule}\leaders\hrule height 0.4pt\hfill}}\vfill
  }}

% Prepare \ScriptRest for slide #1 (zero-padded id) and return nothing.
\newcommand{\PrepareScript}[2]{%
  \renewcommand{\ScriptFile}{build/tex/script-#1.tex}%
  \FitScript
  \global\setbox\ScriptRest\expandafter\ScriptBox\ChosenSize
  \ifScriptOver\typeout{SCRIPT-CONTINUED slide #2}\fi}

\newcommand{\ColumnOverlay}[1]{%
  \begin{tikzpicture}[remember picture,overlay]
    \fill[colbg] ([xshift=-\ColW]current page.north east) rectangle (current page.south east);
    \draw[colrule,line width=0.6pt] ([xshift=-\ColW]current page.north east) -- ([xshift=-\ColW]current page.south east);
    \node[anchor=north west,inner sep=0pt] at ([xshift=-\ColW+5mm,yshift=-11mm]current page.north east) {#1};
  \end{tikzpicture}}

% extra pages when the script is longer than one column
\newcommand{\ContinueScript}[1]{%
  \ifvoid\ScriptRest\else
    \clearpage
    \leavevmode\hbox to 0pt{\ColumnOverlay{\ColumnPage{SCRIPT (CONT.)}{#1}}\hss}%
    {\sffamily\small\color{ink}Slide #1: script continues in the column \par}%
    \ContinueScript{#1}%
  \fi}

\newcommand{\SpeakerPage}[2]{%
  \renewcommand{\CurSlide}{#2}%
  \leavevmode\hbox to 0pt{\ColumnOverlay{\PrepareScript{#1}{#2}\ColumnPage{SCRIPT}{#2}}\hss}%
  \SlideFrame{build/slides/slide-#1.png}\par
  \ifSpeakerHandout\vspace{4mm}{\small\input{build/tex/handout-#1.tex}\par}\fi
  \ContinueScript{#2}%
  \clearpage}
